\documentclass[12pt,a4paper,twoside]{book}%      autres choix : book, report
\usepackage[utf8]{inputenc}%           gestion des accents (source)
\usepackage[T1]{fontenc}%              gestion des accents (PDF)
\usepackage[french]{babel}%          gestion du fran�ais
\usepackage{textcomp}%                 caract�res additionnels
\usepackage{mathtools,amssymb}
\usepackage{amsthm}
\usepackage{xurl}
\usepackage{mathrsfs}
\usepackage{amsmath}% packages de l'AMS + mathtools
\usepackage{lmodern}
\usepackage[bottom=2.5cm,left=2.5cm,right=2.5cm,top=2.5cm]{geometry}
\usepackage[colorlinks=true, linkcolor=blue, citecolor=blue, urlcolor=blue]{hyperref}
\usepackage{fancyhdr}

\usepackage{graphicx}
\usepackage{xcolor}
\pagestyle{fancy}
\fancyhf{}
\theoremstyle{definition}
\newtheorem{dfn}{Définition}[section]
\newtheorem*{pr}{\textbf{Preuve}}
\newtheorem{rem}{\textit{Remarque}}[section]
\newtheorem{thm}[dfn]{\textit{Théorème}}
\newtheorem{prop}[dfn]{\textit{Proposition}}
\newtheorem{cor}[dfn]{\textit{Corollaire}}
\newtheorem{lem}[dfn]{\textit{Lemme}}
\newtheorem{exm}[dfn]{\textit{Exemple}}
\newtheorem{prp}[dfn]{\textbf{Propriété}}
\newtheorem{for}[dfn]{\textit{Formule}}
\numberwithin{equation}{section}
\newtheorem{ex}{Exemple}[subsection]
%\setcounter{page}{1}
\title{ \textbf{Théorème Central Limite et Principe de Déviations Modérées pour les Équations Différentielles Stochastiques de McKean-Vlasov)}}
\begin{document}
	\tableofcontents
	\newpage
	\chapter{Rappel et Préliminaire}
\section{Quelque définition}
\begin{dfn}
Soit $(\Omega,\mathcal{F},\mathbb{P})$ un espace de probabilité et $\{\mathcal{F}_t,t\geq0\}$ une famille de sous-tribus de $\mathcal{F}$. On dit que $\{\mathcal{F}_t,t\geq0\}$ est une famille si c'est une famille croissante, au sens où $\mathcal{F}_s\subset\mathcal{F}_t$ si $s\leq t$.
\end{dfn}
\begin{rem}
	Si $\{\mathcal{F}_t,t\geq0\}$ est une filtration sur un espace de probabilité $(\Omega,\mathcal{F},\mathbb{P})$, on dit que $(\Omega,\{\mathcal{F},t\geq 0\},\mathbb{P})$ est un espace de probabilité filtré.
\end{rem}

\begin{dfn}
	Un processus stochastique $X(t,\omega)$ à valeurs dans un espace mesurable $(E,\mathcal{E})$ est une application de $T\times\Omega$ dans $E$ qui est mesurable par rapport à la mesure du produit $\lambda\cdot\mathbb{P}$ où $\lambda$ est la mesure de Lebesgue sur $T$. 
\end{dfn}
\begin{dfn}
	On dit que le processus $X$ est $(\mathcal{F}_t)-$adapté qi $X_t$ est $\mathcal{F}_t-$mesurable pour tout $t\geq0.$
\end{dfn}	
\begin{dfn}
	Une processus $X_t(\omega)$ sur $[0,T]\times\Omega$ est dite progressivement mesurable par rapport à la filtration $\mathcal{F}_t$ si pour tout $t\in[0,T]$ l'application:
	\[
	X: (s,\omega)\to X_t(\omega)
	\]
	de $[0,t]\times\Omega\to\mathbb{R}$ est $\mathcal{B}[0,t]\otimes\mathcal{F}_t-$mesurable.
\end{dfn}
\begin{prop}
Toute processus stochastique adaptée est progressivement mesurable.
\end{prop}
\begin{pr}
	Dans la livre (Calcul stochastique et modelés de diffusion)
\end{pr} 
\begin{dfn}
	Mouvement brownien
\end{dfn}
\end{document}